\chapter{Bao tải linh kiện}
% full version: chapters/01b_neurontypes.tex

\pencilfig{1}{\pencilart{1}}{Bao tải linh kiện: gần như toàn màu xanh, một ít màu đỏ và một nhúm các loại khác.}

Hãy tưởng tượng bạn được giao một bao tải chứa tám mươi sáu tỷ linh kiện trông giống hệt nhau và được yêu cầu tìm hiểu xem chiếc máy tính được lắp từ chúng như thế nào. Bạn sẽ bắt đầu từ đâu? Một kỹ sư sẽ không ngồi ngắm hình dạng linh kiện. Anh ta sẽ hỏi: linh kiện có bao nhiêu đầu vào, bao nhiêu đầu ra, và nó đưa ra đầu ra cái gì.

Nơ-ron --- tên gọi của loại linh kiện này --- có cấu tạo đơn giản đến bất ngờ. Đầu vào của nó có hàng nghìn: dây dẫn từ các nơ-ron khác đi tới nó. Đầu ra chỉ có một. Nhưng đầu ra này phân nhánh như một cái cây, và cùng một bản sao tín hiệu đi tới hàng nghìn nơi nhận. Hãy hình dung một người cùng lúc nghe nghìn người nói, rồi đáp lại bằng một câu duy nhất, nhưng câu đó lại được thêm một nghìn người khác nghe thấy.

Tín hiệu đó là gì? Một tiếng tách ngắn. Thỉnh thoảng nơ-ron phát ra một xung --- một đột biến điện áp dài khoảng một phần nghìn giây. Mọi tiếng tách của một nơ-ron đều như nhau, như những tiếng gõ của cùng một cái trống. Vì vậy mọi điều nơ-ron có thể nói đều không nằm ở độ lớn của tiếng tách, mà ở chỗ \emph{khi nào} nó tách: dày hơn, thưa hơn, ngay lập tức hay có trễ.

Bên trong, nơ-ron làm phép tính đơn giản nhất có thể nghĩ ra: cộng các tín hiệu đến và, nếu tổng vượt ngưỡng, thì chính nó phát ra tiếng tách. Có đầu vào đẩy tổng lên, có đầu vào kéo tổng xuống.

\section*{Phân loại bao tải thế nào}

Đầu mút của dây dẫn ra, khi vươn tới nơ-ron bên cạnh, không chạm vào nó. Giữa chúng còn một khe hẹp, và tín hiệu đi qua khe nhờ hóa học: dây dẫn nhả ra một nhúm chất đặc biệt --- chất dẫn truyền thần kinh --- và nơ-ron bên cạnh bắt lấy nó. Và đây là điều may mắn: gần như mọi nơ-ron đều nhả cùng một chất ở tất cả các đầu mút dây dẫn của mình. Nghĩa là có thể phân loại nơ-ron theo \emph{chất} mà chúng nhả ra.

Chúng ta sẽ gọi loại nơ-ron bằng bốn chữ cái Latin đầu tiên trong tên tiếng Anh của chất đó. Nơ-ron nhả glutamate là \nt{GLUT}. Nơ-ron nhả axit gamma-aminobutyric là \nt{GABA}. Tiếp theo là \nt{DOPA} (dopamine), \nt{SERO} (serotonin), \nt{NORA} (noradrenaline), \nt{ACET} (acetylcholine) và vài loại nữa.

Khi bao tải đã được phân loại xong, bức tranh hiện ra rất chênh lệch. Cứ mỗi nghìn nơ-ron thì khoảng chín trăm sáu mươi là \nt{GLUT}. Tín hiệu của chúng đẩy nơi nhận tới gần tiếng tách: đó là ``dấu cộng''. Khoảng ba mươi sáu nơ-ron nữa là \nt{GABA}, tín hiệu của chúng đẩy nơi nhận ra xa ngưỡng: đó là ``dấu trừ''. Trong vỏ não chúng nhiều hơn, từ một phần năm đến một phần tư. Còn tất cả các loại khác cộng lại chỉ được vài nơ-ron trên một trăm nghìn. Một giọt nước giữa biển.

\section*{Điện thoại và radio}

Nhưng giọt nước này được tổ chức hoàn toàn khác. Dây dẫn của nơ-ron \nt{GLUT} hay \nt{GABA} đi tới những nơ-ron cụ thể, như một đường dây điện thoại: chỉ người được gọi mới nhận được tín hiệu. Còn những nhóm nhỏ xíu \nt{DOPA}, \nt{SERO}, \nt{NORA} và \nt{ACET} hoạt động như các đài phát thanh. Dây dẫn của chúng tỏa ra khắp những vùng não rộng lớn, chất dẫn truyền thường không đổ vào khe mà tràn ra khoảng không xung quanh, và cùng một thông điệp được hàng triệu tế bào nghe thấy một lúc.

Trên các đường dây điện thoại là dữ liệu: đây là một đường kẻ trên bức tranh, đây là một âm có cao độ nhất định. Trên sóng radio là các thiết lập chung: mọi thứ to nhỏ đến đâu, có nên học không, đã đến giờ ngủ chưa. Máy tính nào cũng vậy: ô nhớ có hàng tỷ, còn thanh ghi điều khiển chỉ một nhúm. Nhưng thiếu chúng thì cỗ máy vẫn không chạy.

\section*{Ý nghĩa do nơi nhận quyết định}

Có một điểm tinh tế. Bản thân chất dẫn truyền không biết nó là ``cộng'' hay ``trừ''. Phân tử là một chiếc chìa khóa. Điều gì xảy ra khi tra chìa vào là do ổ khóa quyết định --- một loại protein đặc biệt ở phía nơi nhận. Cùng một lượng dopamine có thể kích thích nơi nhận này mà lại làm dịu nơi nhận kia, tùy theo chúng lắp ổ khóa nào. Giống như cùng một chương trình radio mà mỗi thính giả hiểu theo một kiểu.

\section*{Tín hiệu ``tốt hơn mong đợi''}

Kênh radio nổi tiếng nhất là dopamine. Nó truyền đi điều gì? Đây là một thí nghiệm. Con vật bất ngờ được cho một giọt nước ép, và các nơ-ron dopamine đáp lại bằng một loạt tiếng tách dồn dập. Sau đó, mỗi lần trước khi cho nước ép, người ta bật một bóng đèn. Con vật học được rằng đèn hứa hẹn nước ép, và giờ loạt tiếng tách xảy ra khi đèn sáng, còn với chính nước ép thì không. Còn nếu sau khi đèn sáng mà không cho nước ép, các nơ-ron dopamine \emph{im bặt} đúng vào lúc chờ đợi.

\pencilfig{101}{%
% three rows: a spike raster of a DOPA neuron; lamp (amber) and juice (blue drop) above the axis
\foreach \row/\y in {1/0,2/-1.05,3/-2.1}{
  \draw[pen] (0,\y) -- (6.2,\y);
  \foreach \s in {0.3,0.75,1.1,2.15,2.6,3.05,3.45,3.9,4.45,4.95,5.4,5.85}{
    \pgfmathtruncatemacro{\gap}{(\row==3)&&(\s>3.6)&&(\s<4.7)}
    \ifnum\gap=0 \draw[pcGray, line width=0.7pt] (\s,\y) -- (\s,\y+0.3);\fi}}
% row 1: juice without a lamp -> burst at the juice
\foreach \s in {4.0,4.08,4.16,4.24,4.33}{\draw[pcAmber, line width=0.8pt] (\s,0) -- (\s,0.45);}
\draw[dotpen=pcBlue, wash=pcBlue] (4.15,0.72) circle (0.09); \draw[pcBlue, line width=0.6pt] (4.07,0.76) -- (4.15,0.92) -- (4.23,0.76);
% row 2: lamp -> burst at the lamp, nothing at the promised juice
\foreach \y in {-1.05,-2.1}{
  \foreach \s in {1.5,1.58,1.66,1.74,1.83}{\draw[pcAmber, line width=0.8pt] (\s,\y) -- (\s,\y+0.45);}
  \draw[dotpen=pcAmber, wash=pcAmber] (1.65,\y+0.72) circle (0.1);
  \foreach \a in {30,90,150}{\draw[pcAmber, line width=0.5pt] ($(1.65,\y+0.72)+(\a:0.14)$) -- ($(1.65,\y+0.72)+(\a:0.24)$);}}
\draw[dotpen=pcBlue, wash=pcBlue] (4.15,-0.33) circle (0.09); \draw[pcBlue, line width=0.6pt] (4.07,-0.29) -- (4.15,-0.13) -- (4.23,-0.29);
% row 3: lamp, no juice -> a gap where the juice was expected
\draw[pcBlue, line width=0.6pt, densely dotted] (4.15,-1.38) circle (0.09); \draw[pcBlue, line width=0.6pt, densely dotted] (4.07,-1.34) -- (4.15,-1.18) -- (4.23,-1.34);
\draw[penlite=pcRed] (3.65,-2.22) -- (3.65,-2.28) -- (4.65,-2.28) -- (4.65,-2.22);
\node[lbl, text=pcRed, anchor=north] at (4.15,-2.3) {khoảng lặng};
\node[lbl, anchor=east, align=right] at (-0.1,0.15) {nước ép\\không có đèn};
\node[lbl, anchor=east, align=right] at (-0.1,-0.9) {đèn,\\rồi nước ép};
\node[lbl, anchor=east, align=right] at (-0.1,-1.95) {đèn,\\không nước ép};
}{Loạt tiếng tách khi có nước ép bất ngờ; khi đèn hứa hẹn nước ép bật sáng; và khoảng lặng đúng lúc nước ép đã hứa không tới.}

Quy tắc giải thích cả ba trường hợp: dopamine không báo ``tốt'' mà báo ``tốt hơn tôi mong đợi''. Nước ép bất ngờ là một bất ngờ dễ chịu. Nước ép đã được hứa thì chẳng có gì bất ngờ. Đã hứa mà không cho là một bất ngờ khó chịu. Đúng đại lượng này được dùng trong những chương trình học từ chính sai lầm của mình. Chúng ta sẽ còn quay lại với nó.

\section*{Điều cần mang theo}

Bộ não là một mạng khổng lồ gồm những linh kiện trông giống hệt nhau. Gần như tất cả là ``cộng'' và ``trừ'' của mạng điện thoại truyền dữ liệu. Chỉ rất ít là các đài phát thanh, chỉnh cả mạng cùng một lúc. Ở chương sau chúng ta sẽ thử xem có thể lắp được gì chỉ từ những ``dấu cộng''.

\fullversion{1}
